La mesure du $\mathrm{pH}$ d'une solution aqueuse d'acide éthano\"ique, à
$\qty{25}{\degreeCelsius}$, a donné $\mathrm{pH}=3,0$.

\begin{enumerate}
    \item Écrire l'équation chimique modélisant la transformation entre l'acide
          éthanoïque et l'eau.
    \item Déterminer l'espèce du couple
          $(\ce{CH3COOH_{(aq)}}\ttslash{}\ce{CH3COO^-_{(aq)}})$ qui prédomine
          dans la solution. Justifier.
    \item Déterminer la valeur du quotient de réaction $Q_{r,\eq}$ à l'état
          d'équilibre du système chimique.
    \item La valeur de $Q_{r,\eq}$ est-elle modifiée si on dilue la solution
          d'acide éthanoïque~? Justifier.
\end{enumerate}

